\chapter{Applications and Reductions}

\section{Purpose of the Applications Chapter}

The preceding chapters define the common URF mechanism:

\[
\text{entropy growth}
+
\text{capacity bound}
+
\text{terminal condition}
\Longrightarrow
\text{refinement depth}.
\]

Applications begin when a domain-specific problem is translated into this
mechanism. A translation is not automatically a proof of the original problem.
It is a reduction surface: it records which objects have been identified, which
bounds have been proved, and which assumptions remain open.

This chapter explains how URF-style reductions appear in computational,
spectral, physical, graph-theoretic, and structural settings.

\section{Reduction Data}

\begin{definition}[Application reduction datum]
An application reduction datum is a tuple
\[
(\mathcal D,X,\mathcal R,H,C,T,\Pi),
\]
where \(\mathcal D\) is the source domain, \(X\) is the translated state space,
\(\mathcal R\) is the admissible refinement class, \(H\) is the entropy or
defect functional, \(C\) is the capacity bound, \(T\) is the terminal set, and
\(\Pi\) is the interpretation map back to the source domain.
\end{definition}

The interpretation map \(\Pi\) is essential. Without it, the URF object may be
internally valid but disconnected from the application it is supposed to
represent.

\begin{definition}[Sound reduction]
A reduction datum is sound if every admissible source-domain process under
study is represented by an admissible refinement trajectory in \(X\), and if
terminal resolution in \(X\) implies the required source-domain conclusion.
\end{definition}

Soundness is the main burden in applications. It is usually harder than the
capacity-depth calculation itself.

\section{Computational Complexity}

In computational applications, the state space \(X\) is often a space of
partial information states: partitions, transcripts, local views, candidate
assignments, proof states, or refinement histories.

\begin{definition}[Computational refinement model]
A computational refinement model consists of a configuration space \(\Omega\),
a refinement state space \(X\), and an admissible update class
\(\mathcal R\) representing the allowed algorithmic steps.
\end{definition}

The entropy \(H\) measures unresolved computational alternatives. For a block
\(B\subseteq \Omega\), a basic entropy is
\[
H(B)=\log |B|.
\]

\begin{theorem}[Computational Depth Transfer]
Suppose a computational problem reduces to a refinement model
\((X,\mathcal R,H,C,T)\). If every admissible update removes at most \(C\)
entropy and every correct solution requires reaching \(T=\{H=0\}\), then every
admissible computation represented by the model has depth at least
\[
\frac{H(x_0)}{C}.
\]
\end{theorem}

\begin{proof}
This is the EntropyDepth lower bound applied to the computational refinement
trajectory.
\end{proof}

\section{SAT-Style Reductions}

For a Boolean search problem on \(n\) variables, the full assignment space has
size \(2^n\). The initial assignment entropy is therefore
\[
\log_2(2^n)=n.
\]

If a local refinement method removes at most \(C\) bits per admissible step,
then terminal isolation of one assignment requires at least \(n/C\) steps.

\begin{remark}[Boundary of SAT-style statements]
This does not prove a SAT lower bound for arbitrary algorithms. It proves a
lower bound only for algorithms already shown to fit the bounded-capacity
refinement model. A theorem about \(P\) versus \(NP\) would require a separate
normalization theorem showing that the relevant algorithms must pass through
such a refinement model.
\end{remark}

\section{Spectral Rigidity}

Spectral applications use operators to measure coercive obstruction.

Let \(L\) be a nonnegative self-adjoint operator with eigenvalues
\[
0=\lambda_0\le \lambda_1\le \lambda_2\le\cdots.
\]

\begin{definition}[Spectral gap]
The first positive spectral gap is a number \(\lambda_1>0\) such that
\[
\langle f,Lf\rangle\ge \lambda_1\|f\|^2
\]
for all \(f\perp \ker L\).
\end{definition}

\begin{theorem}[Spectral Rigidity]
If \(L\) has a positive spectral gap \(\lambda_1>0\), then every
\(f\perp\ker L\) satisfies
\[
\|f\|^2\le \lambda_1^{-1}\langle f,Lf\rangle.
\]
\end{theorem}

\begin{proof}
The spectral gap inequality gives
\[
\langle f,Lf\rangle\ge \lambda_1\|f\|^2.
\]
Dividing by \(\lambda_1>0\) gives the stated estimate.
\end{proof}

This estimate says that deformation away from the kernel has measurable cost.
In URF terms, the spectral gap is a coercivity mechanism.

\section{Physics}

In physical applications, the refinement state may represent accessible
measurements, boundary data, detector outcomes, energy shells, spectral modes,
or coarse-grained macrostates. The capacity bound may come from finite detector
resolution, finite energy, finite channel width, locality, or bounded coupling.

\begin{definition}[Operational physical reduction]
An operational physical reduction represents a physical system by a state
space \(X\), an admissible measurement or update class \(\mathcal R\), an
entropy \(H\), and a terminal condition \(T\) corresponding to the physical
resolution claim.
\end{definition}

\begin{example}[Finite detector capacity]
If a detector has \(B\) distinguishable outcomes per admissible measurement,
then one measurement carries at most \(\log B\) outcome information. If the
unresolved operational entropy is \(H_0\), then at least \(H_0/\log B\)
measurements are required, assuming each measurement is independent in the
required model and terminal resolution requires eliminating that entropy.
\end{example}

A physical application must state the measurement model. Without that model,
the capacity claim is only an analogy.

\section{Case Study: Projected Axial GfE Quasinormal-Mode Certificate}

A concrete cross-repository example is supplied by the Chronos GfE calculation
merged as commit \texttt{2caf5063b85d0fec1084fd18c8e47c7b0da26f96}.  Its
claim is deliberately restricted to the projected, massless, order-reduced
axial effective-field-theory branch with
\[
\ell=2,\qquad n=0,\qquad
\varepsilon=\frac{\beta^2}{M^4}\in[0,10^{-4}].
\]
The source-side projection and order reduction are part of the certified
statement; they are not consequences for the full unreduced trace-log
operator.

\begin{example}[Certified projected axial branch]
The certificate records a canonical axial master potential, validated ingoing
horizon and outgoing infinity endpoint data, and a complex interval-Newton
matching tube.  For every allowed value of \(\varepsilon\), that tube contains
exactly one simple axial \((\ell,n)=(2,0)\) quasinormal-mode root.

At \(\varepsilon=0\), the certified implicit derivative satisfies
\[
\left.\frac{d\Omega_{220}^{(-)}}{d\varepsilon}\right|_{0}
\in
(0.027340131607167324-0.014681561963391043i)
+\overline B(0,0.008950737656596899).
\]
In particular,
\[
0.018389393950570425
\le
\operatorname{Re}
\left.\frac{d\Omega_{220}^{(-)}}{d\varepsilon}\right|_{0}
\le
0.03629086926376422,
\]
so the real part of the first-order frequency correction is strictly positive
throughout the certified disk.
\end{example}

The reduction surface is reproducible through the Chronos generator
\path{tools/gfe/derive_gfe_axial_quadratic_action.sh}, the machine-readable
record
\path{artifacts/chronos/gfe_axial_220_projected_eft_certificate.json}, and its
runtime receipt.  This is an application of verification-first reduction
soundness: the branch, parameter interval, endpoint conditions, root
uniqueness, derivative enclosure, and excluded sectors are all explicit.

\begin{remark}[Exact claim boundary]
This certificate does not establish the polar master equation or polar root,
does not certify axial--polar splitting, and does not cover the full unreduced
trace-log spectrum or its additional massive mode.  Those are independent
mathematical obligations.
\end{remark}

\subsection{Certified Axial Reduction Datum}

The axial certificate admits an explicit defect-form realization of the
application reduction datum. Define the matching defect at \(z=3\) by
\[
\mathfrak D(\Omega,\varepsilon)
=
{\left[
P\,\partial_z\log\Phi_{\mathrm{hor}}-i\Omega
\right]}_{z=3}
-
q_{\mathrm{out}}(3;\Omega,\varepsilon).
\]

\begin{theorem}[Projected axial application datum]
For the projected massless order-reduced axial
\((\ell,n)=(2,0)\) branch, the following tuple is an application reduction
datum:
\[
\left(
\mathcal D_{\mathrm{ax}},
X_{\mathrm{ax}},
\mathcal R_{\mathrm{ax}},
H_{\mathrm{ax}},
C_{\mathrm{ax}},
T_{\mathrm{ax}},
\Pi_{\mathrm{ax}}
\right).
\]
Its components are as follows.

\begin{description}
\item[\(\mathcal D_{\mathrm{ax}}\)]
The projected massless order-reduced axial boundary-value problems with
\[
\varepsilon=\frac{\beta^2}{M^4}\in[0,10^{-4}],
\qquad
\ell=2,
\qquad
n=0.
\]

\item[\(X_{\mathrm{ax}}\)]
Affine complex enclosures for \(\varepsilon\), \(\Omega\), the horizon
logarithmic derivative, and the outgoing infinity logarithmic derivative at
the common matching point \(z=3\).

\item[\(\mathcal R_{\mathrm{ax}}\)]
Validated outward Riccati propagation from the horizon, validated inward
Riccati propagation from the outgoing infinity endpoint, and affine complex
interval-Newton contraction of the matching defect.

\item[\(H_{\mathrm{ax}}\)]
The nonnegative defect functional
\[
H_{\mathrm{ax}}(\Omega,\varepsilon)
=
\left|\mathfrak D(\Omega,\varepsilon)\right|.
\]

\item[\(C_{\mathrm{ax}}\)]
The certified Newton enclosure capacity
\[
r_{\mathrm N}=\frac{1}{150000},
\]
together with the verified control inequality
\[
\rho_{\mathrm N}
=
\frac{\text{analytic error}}
     {\text{derivative lower bound}}
<
r_{\mathrm N}.
\]

\item[\(T_{\mathrm{ax}}\)]
The family of Newton tubes in which
\(\mathfrak D(\Omega,\varepsilon)=0\) has exactly one simple root for every
\(\varepsilon\in[0,10^{-4}]\).

\item[\(\Pi_{\mathrm{ax}}\)]
The interpretation map sending each certified terminal tube to its unique
projected axial frequency
\(\Omega_{220}^{(-)}(\varepsilon)\), and at \(\varepsilon=0\) sending the
certified tangent data to
\[
\left.
\frac{d\Omega_{220}^{(-)}}{d\varepsilon}
\right|_0
=
-\frac{D_{\varepsilon}}{D_{\Omega}}.
\]
\end{description}
\end{theorem}

\begin{proof}
The source domain and excluded sectors are fixed by the machine-readable
certificate. The two endpoint propagations and interval-Newton calculation
supply \(X_{\mathrm{ax}}\) and \(\mathcal R_{\mathrm{ax}}\). The matching
quantity in the generator defines \(H_{\mathrm{ax}}\). The strict internal
Newton-image inequality supplies \(C_{\mathrm{ax}}\), while the certified
unique-simple-root statement supplies \(T_{\mathrm{ax}}\). The frequency and
implicit-derivative records supply \(\Pi_{\mathrm{ax}}\).
\end{proof}

\begin{remark}[Defect-form boundary]
This is a soundness and certification datum, not an entropy-depth lower-bound
claim. Here \(C_{\mathrm{ax}}\) is a validated enclosure capacity rather than
a per-step information-capacity constant.
\end{remark}

\subsection{Constraint-Preserving Derivative Reductions}

\begin{definition}[Algebraic constraint at fixed perturbative order]
Let \(\widehat S[\widehat C,\widehat q]\) be a transformed action truncated at a
fixed perturbative order. The field \(\widehat C\) remains an algebraic
constraint at that order if its Euler--Lagrange expression
\(E_{\widehat C}(\widehat S)\) has differential order zero in \(\widehat C\).
Derivatives of the remaining fields \(\widehat q\) may occur.
\end{definition}

\begin{theorem}[Derivative obstruction to algebraic constraint preservation]
Suppose a proposed perturbative reduction requires \(\widehat C\) to remain an
algebraic constraint. If the simplified Euler--Lagrange expression contains
\[
a_{\alpha}\,\partial^{\alpha}\widehat C,
\qquad
|\alpha|\geq 1,
\]
with \(a_{\alpha}\) not identically zero on the claimed domain, then the
reduction is not algebraic-constraint preserving on that domain. In
particular, \(\widehat C\) cannot be eliminated there by an algebraic
constraint equation.
\end{theorem}

\begin{proof}
An algebraic constraint has differential order zero in \(\widehat C\).
A nonzero derivative jet \(\partial^{\alpha}\widehat C\), with
\(|\alpha|\geq 1\), gives positive differential order and contradicts the
definition.
\end{proof}

\subsection{Polar Constraint-Preservation Obstruction}

The polar continuation sharpens the preceding boundary without closing it.
Chronos first certifies the Schwarzschild polar
\((\ell,n)=(2,0)\) root at \(\varepsilon=0\) by an exact Darboux transfer.
At relative order \(O(\beta^2)\), however, the unreduced even-parity
quadratic-action source has a rank-two fourth-time principal sector.  An
explicit perturbative derivative substitution removes that temporal sector
modulo a total derivative before any lapse or shift constraint is solved.

The transformed lapse equation passes the temporal-order test: its highest
remaining time derivative is second order, and no third- or higher-time
derivative survives.  This does not preserve the usual algebraic lapse
constraint, because the static \(H_0\)-only radial sector contains
\[
\mathcal L^{(2)}_{\beta^2,H_0,\mathrm{rad}}
=
\frac{4M{(r-2M)}^2}{3r^3}{\bigl(H_{0,rr}\bigr)}^2.
\]
Consequently the \(H_0\) Euler equation contains the nonzero exterior
principal term
\[
\frac{8M{(r-2M)}^2}{3r^3}H_{0,rrrr},
\qquad M>0,\quad r>2M.
\]

\begin{theorem}[Polar lapse non-algebraicity]
On the Schwarzschild exterior \(M>0\), \(r>2M\), the current transformed polar
branch is not algebraic-constraint preserving for \(H_0\). Hence the standard
algebraic lapse elimination is unavailable on that domain.
\end{theorem}

\begin{proof}
The coefficient of \(H_{0,rrrr}\) is
\[
\frac{8M{(r-2M)}^2}{3r^3},
\]
which is strictly positive for \(M>0\) and \(r>2M\). The derivative-obstruction
theorem therefore applies with \(\widehat C=H_0\).
\end{proof}

\begin{remark}[Constraint-preservation boundary]
The time-only order-reduction map changes \(H_2\) and \(K\), not \(H_0\), and
therefore cannot cancel this \(H_0\)-only radial term.  The current projected
polar branch does not admit the standard algebraic elimination of the lapse.
No \(H_0\) solution, \(H_1\) elimination, polar master equation,
\(\varepsilon>0\) polar root continuation, or axial--polar splitting follows
from this result.
\end{remark}

The executable obstruction is recorded in Chronos commit
\texttt{c99961425877c9d49117b545f5eb22ef6a5ddb43}, with the verifier
\path{tests/test_gfe_polar_h0_radial_derivative_obstruction.py}.  This case
study illustrates a reduction-soundness failure detected before introducing a
master variable: removing the highest time derivatives is insufficient when a
would-be constraint field acquires an independent higher-radial equation.


\subsection{Unreduced Six-State Invariant-Graph Boundary}

The unreduced axial equations were next tested directly in the dimensionless
six-state derivative-jet variables
\[
q=\frac{\beta}{M^2},
\qquad
x=\frac rM,
\qquad
\Omega=M\omega,
\qquad
Z=(H_0,H_{0,x},H_{0,xx},H_1,H_{1,x},H_{1,xx}).
\]
For \(0<q\le10^{-2}\) and \(x>2\), the exact coefficientwise principal solve
has determinant
\[
\Delta
=
\frac{
4q^2(-16q+5x^3)(8q+5x^3){(x-2)}^3
}{x^{11}},
\]
which is strictly positive throughout the open exterior.

Let \(P_0\) be the leading GR massless projector and
\(Q_0=I-P_0\).  Writing the complementary graph component as \(y=p\,u\),
the exact invariance equation is
\[
R(q,p)
=
p_x-A_{QP}-A_{QQ}p+pA_{PP}+pA_{PQ}p.
\]
Analytic coefficient matching gives
\[
p(q,x)=q^2P_2(x)+O(q^3)
\]
and forces
\[
{(P_2)}_{\delta H_{1,xx},H_0}
=
\frac{
4i\Omega(108x^3-549x^2+653x-64)
}{
9x^5{(x-2)}^4
}.
\]

The previously certified ingoing-EF multiplier is
\[
W=\operatorname{diag}(1,1,f,f^3,f^2,f),
\qquad
f=\frac{x-2}{x}.
\]
Under componentwise application of that multiplier to the derivative-jet
six-state graph,
\[
\lim_{x\downarrow2}
{(x-2)}^3 f
{(P_2)}_{\delta H_{1,xx},H_0}
=
-\frac{5i\Omega}{8},
\]
which is nonzero on the axial \((2,0)\) frequency tube.  The forced first graph
tangent is therefore not bounded in the proposed horizon-weighted graph
space, so the requested contraction fails at the horizon before an infinity
contraction estimate can close.

\begin{remark}[Exact continuation boundary]
This test does not produce a unique analytic unreduced massless-branch
continuation theorem on \(0\le q\le10^{-2}\).  Consequently the existing
projected axial \((2,0)\) quasinormal-mode, tail, and decay results are not
transferred to the unreduced six-state system.

The obstruction is specifically to componentwise use of the existing EF
multiplier on the derivative-jet representation.  It does not prove
nonexistence of an invariant massless branch under a separately derived
state-conjugate six-state weight.
\end{remark}

\section{Graph Structures}

Graph applications often use local neighborhoods, partition refinements, color
refinements, or spectral graph operators.

\begin{definition}[Graph refinement datum]
A graph refinement datum consists of a graph family \(G_n\), a local
observation rule, a partition or label state space, a refinement rule, and a
terminal separation condition.
\end{definition}

If bounded-radius local views leave \(N_n\) globally compatible alternatives,
then the unresolved graph entropy is
\[
H_n=\log N_n.
\]
If each refinement round removes at most \(C\) units of this entropy, then the
terminal depth is at least \(H_n/C\).

\begin{remark}
This is not a graph-isomorphism lower bound unless the refinement model is
proved to capture the graph-isomorphism procedure under study.
\end{remark}

\section{Expanders and Rigidity}

Expanders provide a useful contrast. A graph family with a uniform spectral
gap has strong global connectivity from local incidence data. In URF language,
a gap can act as a coercive control preventing low-cost deformation.

However, expansion alone is not a complete URF lower bound. A full application
must still specify entropy, capacity, terminal resolution, and the admissible
refinement process.

\begin{theorem}[Gap-to-Coercivity Template]
Let \(L_n\) be graph Laplacians with a uniform spectral gap
\[
\lambda_1(L_n)\ge \lambda>0.
\]
Then for every \(f\perp\ker L_n\),
\[
\|f\|^2\le \lambda^{-1}\langle f,{L_n}f\rangle.
\]
\end{theorem}

\begin{proof}
Apply the spectral gap inequality with \(\lambda_1(L_n)\ge\lambda\).
\end{proof}

\section{Mathematical Structures}

In algebraic, geometric, or categorical settings, an application may use
defect functionals instead of entropy directly. A defect functional
\(F:X\to\mathbb R_{\ge 0}\) becomes URF-relevant when it controls or is
controlled by entropy.

\begin{definition}[Entropy-compatible defect]
A defect functional \(F\) is entropy-compatible if there is a monotone function
\(\psi\) such that
\[
H(x)\le \psi(F(x))
\]
or
\[
F(x)\le \psi(H(x)),
\]
depending on the direction needed by the argument.
\end{definition}

The direction matters. A lower bound on depth requires a quantity that cannot
be eliminated faster than the capacity bound allows.

\section{Reduction Soundness}

The central risk in applications is an unsound translation.

\begin{theorem}[Reduction Soundness Requirement]
A URF lower bound transfers to a source-domain problem only if every
source-domain process in the claimed class is represented by an admissible
refinement trajectory and every source-domain success state maps into the
declared terminal set.
\end{theorem}

\begin{proof}
If a source-domain process is not represented by an admissible trajectory, the
URF lower bound does not apply to it. If source-domain success does not require
entry into the terminal set, the terminal-depth lower bound does not imply a
source-domain lower bound. Therefore both representation and terminal
compatibility are necessary for transfer.
\end{proof}

\section{Application Checklist}

A complete application should provide:

\begin{center}
\begin{tabular}{ll}
\toprule
Field & Required content \\
\midrule
Source domain & problem or system being translated \\
State space & URF representation of source states \\
Admissible process & refinement class modeling source operations \\
Entropy or defect & unresolved quantity being measured \\
Capacity bound & per-step information limit \\
Coercivity & rigidity estimate, if needed \\
Terminal set & translated success condition \\
Soundness bridge & proof that source processes enter the model \\
Boundary & explicit non-claims \\
\bottomrule
\end{tabular}
\end{center}

\section{Boundary}

This chapter gives application templates and reduction requirements. It does
not prove any source-domain lower bound without a soundness bridge, does not
prove SAT hardness, does not prove graph-isomorphism lower bounds, does not
establish any physical theory, and does not prove unrestricted Chronos-RR,
unrestricted H4.1/FGL, P vs NP, or any Clay problem.
